%! AP Statistics — Unit 7, Lesson 7.2 — Guided Notes (Answer Key)
\documentclass[10pt]{article}
\usepackage{apstats-article}
\usepackage{apstats-key}
\newcommand{\TallMath}[1]{$\displaystyle #1\rule[-1.4em]{0pt}{3.2em}$}

\begin{document}

\pageheader{Unit 7, Lesson 7.2}{Guided Notes --- Answer Key}
\namedateperiod

\begin{objectivebox}
By the end of this lesson, I will be able to:
\begin{itemize}
    \item[$\square$] \ansline{Describe the $t$-distribution and how it differs from $N(0,1)$.}
    \item[$\square$] \ansline{Verify the conditions for a one-sample $t$-interval.}
    \item[$\square$] \ansline{Calculate a one-sample $t$-interval: $\bar x \pm t^* \cdot s/\sqrt{n}$.}
    \item[$\square$] \ansline{Apply the one-sample $t$-interval to matched-pairs differences.}
\end{itemize}
\end{objectivebox}

\begin{vocabbox}
\noindent\textbf{\textcolor{navy}{$t$-distribution:}}\\[1pt]%
\ansline{Symmetric, bell-shaped curve with heavier tails than $N(0,1)$; used when $s$}\\[1pt]\ansline{replaces an unknown $\sigma$. Indexed by degrees of freedom.}\\[3pt]
\noindent\textbf{\textcolor{navy}{Degrees of freedom ($df$):}}\\[1pt]%
\ansline{The parameter that picks out one $t$-curve. One-sample $t$: $df = n - 1$.}\\[1pt]\ansline{}\\[3pt]
\noindent\textbf{\textcolor{navy}{Standard error of $\bar x$ ($SE_{\bar x}$):}}\\[1pt]%
\ansline{$SE_{\bar x} = s/\sqrt{n}$; estimates the SD of the sampling distribution of $\bar x$.}\\[1pt]\ansline{}\\[3pt]
\noindent\textbf{\textcolor{navy}{Critical value $t^*$:}}\\[1pt]%
\ansline{Value with area $C$ between $-t^*$ and $t^*$ under the $t$-curve;}\\[1pt]\ansline{found from the $t$-table (or invT) with $df = n - 1$.}\\[3pt]
\noindent\textbf{\textcolor{navy}{One-sample $t$-interval:}}\\[1pt]%
\ansline{$\bar x \pm t^* \cdot s/\sqrt{n}$; a CI for a population mean $\mu$ when $\sigma$ is unknown.}\\[1pt]\ansline{}\\[3pt]
\noindent\textbf{\textcolor{navy}{Matched-pairs design:}}\\[1pt]%
\ansline{Each observation is paired with another (same subject or matched subjects);}\\[1pt]\ansline{analyze the differences with a one-sample $t$-procedure.}\\[3pt]
\end{vocabbox}

\begin{hookbox}
A nutritionist records daily sodium intake for a random sample of 20 college students:
$\bar x = 2840$ mg, $s = 480$ mg. The recommended maximum is 2300 mg.

Why can't we use a $z$-interval to estimate the true mean sodium intake? \ansline{$\sigma$ is unknown.}\\
\ansline{Using $s$ in place of $\sigma$ makes the standardized statistic follow a $t$, not $z$.}
\end{hookbox}

\begin{notesbox}{The $t$-Distribution (VAR-7.A)}
When we substitute $s$ for $\sigma$, the statistic $t = \dfrac{\bar x - \mu}{s/\sqrt{n}}$
follows a \ans{$t$}-distribution.

\vspace{0.1in}
\textbf{Properties:}
\begin{itemize}
    \item Bell-shaped, \ans{symmetric}, centered at \ans{0}.
    \item Heavier tails than the standard normal (because $s$ adds variability).
    \item Indexed by \textbf{degrees of freedom}: $df = n - $ \ans{1}.
    \item As $df$ \ans{increases}, the $t$-distribution approaches $N(0, 1)$.
\end{itemize}

\textbf{Finding $t^*$:} For 95\% CI with $n = 20$: $df = $ \ans{19}, $t^* = $ \ans{2.093}.
\end{notesbox}

\begin{notesbox}{Conditions for a One-Sample $t$-Interval (UNC-4.P)}
\begin{enumerate}
    \item \textbf{Random:} Data come from a \ans{random sample} or randomized experiment.
    \item \textbf{10\%:} If sampling w/o replacement: $n \le $ \ans{10}$\% \cdot N$.
    \item \textbf{Large Sample:}
          \begin{itemize}
              \item $n \ge $ \ans{30}: CLT applies $\to$ condition met.
              \item $n < $ \ans{30}: check that the data show no strong
                    \ans{skewness} or \ans{outliers}.
          \end{itemize}
\end{enumerate}
\end{notesbox}

\begin{notesbox}{The One-Sample $t$-Interval Formula (UNC-4.Q, UNC-4.R)}
\renewcommand{\arraystretch}{2}
\begin{tabularx}{\linewidth}{lX}
    Point estimate & $\bar x$ \\
    Standard error & $SE_{\bar x} = \dfrac{s}{\sqrt{n}}$ \\
    Margin of error & $ME = t^* \cdot \dfrac{s}{\sqrt{n}}$ \quad ($t^*$ from table, $df = n - 1$) \\
    \textbf{Interval} & $\bar x \pm t^* \cdot \dfrac{s}{\sqrt{n}}$ \\
\end{tabularx}

\vspace{0.1in}
\textbf{Sodium example:} $n = 20$, $\bar x = 2840$, $s = 480$, 95\% CI.

$df = $ \ans{19}, \quad $t^* = $ \ans{2.093}, \quad $SE = 480/\sqrt{20} = $ \ans{107.3} mg

Interval: $2840 \pm $ \ans{2.093} $\cdot$ \ans{107.3} $= 2840 \pm$ \ans{224.6}
$= ($ \ans{2615.4}$,$ \ans{3064.6}$)$ mg
\end{notesbox}

\begin{notesbox}{The Four-Step Procedure (PCCF)}
\begin{enumerate}
    \item \textbf{Parameter:} Let $\mu = $ the true mean \ans{[quantity]} for \ans{[population]}.
    \item \textbf{Conditions:} State and verify all three (Random, 10\%, Large Sample).
    \item \textbf{Calculate:} $\bar x \pm t^* \cdot s/\sqrt{n}$
    \item \textbf{Conclude:} ``We are \ans{$C$}\% confident that the interval
          $(\rule{1.5cm}{0.4pt},\ \rule{1.5cm}{0.4pt})$ captures the true mean
          \ans{[quantity]} for \ans{[population]}.''
\end{enumerate}
\end{notesbox}

\begin{notesbox}{Matched-Pairs Design (UNC-4.O.3)}
When data come in pairs (before/after, two treatments on same subject):
\begin{enumerate}
    \item Compute differences $d_i = x_{i,1} - x_{i,2}$ for each pair.
          (State the order of subtraction: \ansline{e.g., after $-$ before; use it consistently.})
    \item Find $\bar d = $ \ans{mean of the $d_i$} and $s_d = $ \ans{SD of the $d_i$}.
    \item Apply the one-sample $t$-interval to the differences:
          $\bar d \pm t^* \cdot \dfrac{s_d}{\sqrt{n}}$, \quad $df = n - 1$,
          where $n = $ number of \ans{pairs}.
\end{enumerate}
\end{notesbox}

\end{document}
